\documentclass{article}

\textwidth=6.5in
\textheight=8.5in
\voffset=-54pt
\oddsidemargin=0in
\evensidemargin=0in
\setlength{\parskip}{8pt}
\setlength{\parindent}{0pt}

\usepackage{amsmath, amssymb}
\usepackage{ifthen}

\usepackage[inline]{enumitem}
\setlist{topsep=0pt, parsep=8pt}

\usepackage[rgb,table]{xcolor}\convertcolorsUfalse
\usepackage{graphicx}

% change hyperref options
\PassOptionsToPackage{hyphens}{url}
\usepackage[bookmarks,colorlinks,linkcolor=black,citecolor=blue,pdfstartview=FitH,urlcolor=blue]{hyperref}
\hypersetup{pdfpagemode=UseNone}
\usepackage{bookmark}

\usepackage{graphicx}
\usepackage{multicol}
\usepackage{framed}
\usepackage{array}

\usepackage{icomma}

\usepackage{marvosym}

\usepackage{pifont}

\usepackage{soul}

\usepackage{tikz}
\usetikzlibrary{
  matrix,%
  % enables the snake paths
  decorations.pathmorphing,%
  % enables bigger arrows
  decorations.markings,%
  decorations.pathreplacing,%
  decorations.text,%
  calc,%
  patterns,%
  shapes.geometric,%
  arrows,
  decorations.shapes,
  positioning,
  plotmarks,
  intersections
}
\tikzset{>=latex} % sets default arrow type in tikz
\tikzset{font=\normalsize}
\tikzset{mark size=1.5pt, mark options=thin}
\tikzset{pin distance=4pt,
  every pin edge/.style={<-, thin, shorten <= -2pt}}

\interfootnotelinepenalty=10000
\renewcommand{\thefootnote}{(\arabic{footnote})}

\usepackage{multirow, bigdelim}

% change to \usepackage[sols]{handout} to see solutions
\usepackage[]{handout}

\providecommand{\check}{}
\renewcommand{\check}{\ifmmode{\text{ \ding{51} }}\else{\ding{51} }\fi}

\newcommand\iscurrentenumi[1]{%
  \ifthenelse{\equal{#1}{\theenumi}}{}{#1}}
\setenumerate[1]{ref=\#\arabic*}
\setenumerate[2]{ref=\protect\iscurrentenumi{\theenumi}(\alph*)}
\newcommand\enumiiprefix[2]{%
  \ifthenelse{{\equal{#1}{\theenumi}}\AND{\equal{#2}{(\alph{enumii})}}}{}{}%
  \ifthenelse{{\equal{#1}{\theenumi}}\AND{\NOT{\equal{#2}{(\alph{enumii})}}}}{#2}{}%
  \ifthenelse{\equal{#1}{\theenumi}}{}{#1#2}%
}
\setenumerate[3]{ref=\protect\enumiiprefix{\theenumi}{(\alph{enumii})}\roman*}

\newlength{\tipmargin}
\makeatletter

\renewenvironment{framed}{
  \setlength{\tipmargin}{\textwidth-\linewidth}
  \advance\@totalleftmargin-\tipmargin
  \def\FrameCommand{\hspace{\tipmargin}\fbox}%
  \advance\linewidth\tipmargin
  \MakeFramed {\advance\hsize-\width \FrameRestore}%
}
{\endMakeFramed}

\makeatother

\definecolor{purple}{rgb}{0.5,0,1.0}

\usepackage{fancyhdr}
\lhead{\textcolor{white}{This content is protected and may not be shared, uploaded, or distributed.}}
\rhead{}
\rfoot{\vspace*{\fill}\footnotesize\textcopyright{} \emph{Harvard Math Ma}}
\renewcommand{\headrulewidth}{0pt}
\pagestyle{fancy}
\begin{document}

\htitle{32. Solving Equations with Logarithms and Exponents}

\begin{problemonly}
  \begin{framed}

You should be able to:
    \begin{itemize}[topsep=0pt,partopsep=0pt,parsep=8pt,itemsep=0pt]
    \item Solve equations involving exponentials and logarithms.

    \end{itemize}
  \end{framed}
\end{problemonly}

\begin{enumerate}
\item  \label{exponentials-demographics}

  \note{This is meant to be motivation for why we might want to solve equations involving exponentials.}

  \begin{problem}[\vfill]
    A country's government is doing some demographic research. They estimate that, $t$ years from now, the number of working-age adults will be $7 (0.98)^t$ million, and the number of retirement-age adults will be $2(1.03)^t$ million. When the number of working-age adults falls to 3 times the number of retirement-age adults, they will need to put in additional measures to finance their health care system, so they would like to predict when this will happen. What equation should they solve to find this out?
  \end{problem}

  \begin{solution}
    We want to determine the value of $t$ that
    satisfies the following equation.
    \[7(0.98)^t = 3\cdot 2(1.03)^t\]
  \end{solution}

\end{enumerate}
In each problem, solve for $x$. 
\begin{enumerate}[resume]
\item 

  \note{Here, I want to emphasize a couple of things:
    \begin{itemize}[nosep]
    \item If an equation is linear in $x$, we can combine the multiples of $x$ on one side of the equation and the constant terms on the other to solve for $x$.
    \item Otherwise, we want to rewrite the equation as something = 0 and then either factor or use the quadratic formula. (This is also a chance for us to remember what the quadratic formula says.)  
    \end{itemize}
  }

  \emph{Warm up.}

  \probonly{\begin{multicols}{2}}
    \begin{enumerate}
    \item \label{solve-linear-warmup}
      \begin{problem}
        $6 (x + 3) = 7 (2x - 1)$
      \end{problem}

      \begin{solution}
        Since this equation is linear in $x$,
        we can solve for $x$ by isolating $x$ on one side
        of the equation.
        \begin{align*}
          6(x+3) &= 7(2x-1) \\
          6x + 18 &= 14x - 7 \\
          18 + 7 &= 14x - 6x \\
          25 &= 8x\\
          x &= \boxed{\frac{25}{8}}
        \end{align*}
      \end{solution}

    \item
      \begin{problem}
        $2x^2 + 1 = -5x$
      \end{problem}
      
      \begin{solution}
        Since this equation is quadratic in $x$, we can solve
        for $x$ by collecting all non-zero terms on one side
        and then factoring or using the quadratic formula.
        \begin{align*}
          2x^2+1 &= -5x \\
          2x^2 + 5x + 1 &= 0 \\
          \intertext{Using the quadratic formula:}
          x &= \frac{-5 \pm \sqrt{5^2-4(2)(1)}}{2(2)} = 
          \boxed{\frac{-5 \pm \sqrt{17}}{4}}
        \end{align*}
      \end{solution}

    \end{enumerate}
  \probonly{\end{multicols}}

  \pspace{\vfill}

\item  \label{exp-log-basic}

  \note{The last two parts of this give us a chance to talk about domain / range before things get more complicated.}

  \emph{3 basic types of equations.} 

  \begin{enumerate}
  \item \label{exp-log-basic-var-in-exponent}
    \begin{problem}[\stretch{0.2}]
      $7^x = 12$
    \end{problem}

    \begin{solution}
      We are looking for the exponent of the power of $7$ that equals 
      $12$. That would be
      $x=\boxed{\log_7 12}$.
    \end{solution}

  \item
    \begin{problem}[\stretch{0.2}]
      $x^7 = 12$
    \end{problem}

    \begin{solution}
      We are looking for the base whose $7$th power equals $12$.
      That would be $x=\boxed{\sqrt[7]{12}}$.
    \end{solution}

  \item
    \begin{problem}[\stretch{0.2}]
      $\log_7 x = 12$
    \end{problem}

    \begin{solution}
      We are looking for the power of $7$ whose exponent equals $12$.
      That would be $x=\boxed{7^{12}}$.
    \end{solution}

  \item

    \note{Here, I want to point out that we can't take a log of a negative number.}

    \begin{problem}[\stretch{0.2}]
      $7^x = -12$
    \end{problem}

    \begin{solution}
      We are looking for the exponent of the power of $7$ that equals
      $-12$. However, there is no power of $7$ that equals $-12$. 
      In other words, $-12$ is not
      in the range of $7^x$. Therefore this equation has
      \fbox{no solution}.
    \end{solution}

  \item
    \begin{problem}[\nstretch{0.2}]
      $\log_7 x = -12$
    \end{problem}

    \begin{solution}
      We are looking for the power of $7$ whose exponent equals $-12$.
      This is $x=\boxed{7^{-12}}$.
    \end{solution}

  \end{enumerate}

\item 

  \note{The second method here is meant to lead into the strategy for the next problem.}

  \begin{problem}
    Now, we'll solve $7^{3x + 1} = 5 \cdot 7^{x + 3}$.
  \end{problem}

  {
    \setlength{\columnseprule}{0.4pt}
    \begin{multicols}{2}
      \textbf{Method 1:} Start by dividing both sides by $7^{x + 3}$ to put the powers of 7 together:

      \begin{solution}
        \begin{align*}
          7^{3x + 1} &= 5 \cdot 7^{x + 3} \\
          \frac{7^{3x+1}}{7^{x+3}} &= 5 \\
          7^{2x-2} &= 5 \\
          2x -2 &= \log_{7}(5) \\
          2x &= \log_7(5)+2 \\
          x &= \boxed{\tfrac{1}{2}\log_7(5) + 1}
        \end{align*}
      \end{solution}

      \pspacestar{2.5in}

      \columnbreak

      \textbf{Method 2:} Since $x$ is in the exponent, start by taking an appropriate log of both sides:

      \pspacestar{2.5in}

      \begin{solution}
        \begin{align*}
          7^{3x + 1} &= 5 \cdot 7^{x + 3} \\
          \log_7(7^{3x + 1}) &= \log_7(5 \cdot 7^{x + 3}) \\
          \log_7(7^{3x + 1}) &= \log_7(5) + \log_7(7^{x+3})\\
          3x+1 &= \log_7(5) + x+3\\
          2x &= \log_7(5)+2 \\
          x &= \boxed{\tfrac{1}{2}\log_7(5) + 1}
        \end{align*}
      \end{solution}
    \end{multicols}
  }

\item 

  \note{These involve multiple bases. In this case, it doesn't really matter what base logarithm we use; our default choice is the natural log.}

  \begin{enumerate}
  \item \label{exponential-equation-two-bases-1}

    \note{Once we take ln of both sides, we end up with an equation that looks scary to students. I encourage them to identify what things are just numbers ($\ln 4$ and $\ln 5$) and imagine replacing them with easier numbers. Here, you exactly get something like \ref{solve-linear-warmup}.} 

    \begin{problem}[\vfill]
      $4^{x + 3} = 5^{2x - 1}$
    \end{problem}

    \begin{solution}
      Let's take the natural log of both sides:
      \begin{align*}
        \ln \left( 4^{x + 3} \right) &= \ln \left( 5^{2x - 1} \right) \\
        (x + 3) \ln 4 &= (2x - 1) \ln 5 \\
        (\ln 4) x + 3 \ln 4 &= (2 \ln 5) x - \ln 5 \\
        3 \ln 4 + \ln 5 &= (2 \ln 5 - \ln 4) x \\
        \boxed{\frac{3 \ln 4 + \ln 5}{2 \ln 5 - \ln 4}} &= x
      \end{align*}
    \end{solution}

  \item

    \note{One thing I'll talk about is whether we can directly take $\ln$ of both sides. I want them to recognize both that we can't simplify $\ln$ of a difference and that $\ln(0)$ is undefined.}

    \begin{problem}[\vfill\newpage]
      $2^{x^2} - 5 e^x = 0$ (How is $2^{x^2}$ different from $(2^x)^2$? What does each one mean if $x = 10$?)
    \end{problem}

    \begin{solution}
      In order to use the same strategy of taking the natural log
      of both sides, we need to first rewrite the equation
      so that both sides consist of a single term.
      \begin{align*}
        2^{x^2} - 5 e^x &= 0 \\
        2^{x^2} &= 5e^x \\
        \ln(2^{x^2}) &= \ln(5e^x) \\
        \ln(2^{x^2}) &= \ln5 + \ln(e^x) \\
        x^2 \ln2 &= \ln5 + x
      \end{align*}
      This is an equation that is quadratic in $x$,
      so we should collect all non-zero terms on one side.
      \[(\ln2)x^2 - x - (\ln5) = 0\]
      Applying the quadratic formula:
      \begin{align*}
        x &= \frac{-(-1) \pm \sqrt{(-1)^2 - 4(\ln2)(-\ln5)}}{2(\ln2)}\\
        x &= \boxed{\frac{1\pm\sqrt{1+2(\ln2)(\ln5)}}{2\ln2}}
      \end{align*}
    \end{solution}

  \item

    \note{This doesn't have anything new, so I'll skip it if we're short on time.}

    \begin{problem}[\vfill]
      Solve the equation from \ref{exponentials-demographics}. 
    \end{problem}

    \begin{solution}
      We can solve the equation by taking the natural log of both sides.
      \begin{align*}
        7(0.98)^t &= 3\cdot 2(1.03)^t \\
        \ln(7(0.98)^t) &= \ln(6(1.03)^t) \\
        \ln 7 + t \ln 0.98 &= \ln 6 + t\ln1.03 \\
        \ln 7 - \ln 6 &= (\ln 1.03)t - \ln(0.98) t \\
        \ln 7 - \ln 6 &= t(\ln1.03 - \ln0.98) \\
        t &= \boxed{\frac{\ln 7 - \ln 6}{\ln 1.03 - \ln0.98}}
      \end{align*}
    \end{solution}
  \end{enumerate}

\item  \label{equations-with-multiple-logs}
\begin{problemonly}
    {\LARGE\Pointinghand} Tips for successful equation solving:
    \begin{itemize}[nosep]
    \item Remind yourself what you're doing in each step. Are you taking the natural log of both sides? Raising $5$ to both sides? Simplifying one side using a log rule?
    \item Work with equations! That is, every step you write down should have an equals sign in it.
    \item Remember that, if you're doing something to one side of an equation (other than just rewriting it), you have to do the same thing to the other side.
    \item Take it one step at a time! Most mistakes happen when you try to do multiple things at once (like taking a log \textit{and} rewriting using log rules). It never hurts to write out more steps.
    \item When solving an equation that has logs in it, it's easy to introduce extraneous solutions. So, at the end, be sure to check that each solution is actually in the domain of the equation.
    \item If you end up with a simple answer, take an extra second to plug it in to the original equation and check whether it works.
    \end{itemize}
  \end{problemonly}
  \begin{enumerate}
  \item

    \note{Here, there are two good strategies. One is to immediately raise $e$ to both sides; another is to first use log rules to rewrite the right side as $\ln(\text{something})$.}

    \begin{problem}[\vfill]
      $\ln(x + 5) = 2 \ln (x + 3)$
    \end{problem}

    \begin{solution}
      Let's start by raising $e$ to both sides:
      \begin{align*}
        e^{\ln (x + 5)} &= e^{2 \ln(x + 3)} \\
        x + 5 &= [e^{\ln (x + 3)}]^2 \\
        x + 5 &= (x + 3)^2 \\
        \intertext{This is an equation that is quadratic in $x$.}
        x + 5 &= x^2 + 6x + 9 \\
        0 &= x^2 + 5x + 4 \\
        0 &= (x + 1)(x + 4) 
      \end{align*}
      So, $x = -1$ and $x = -4$. However, $x = -4$ is not actually a solution because it is not in the domain of the original equation (the original equation only makes sense when $x + 3$ and $x + 5$ are both positive). We say that $x=-4$ is an \textit{extraneous
      solution.} The only solution is $x = \boxed{-1}$.
    \end{solution}

    \note{This will hopefully help students realize that one of the solutions they found isn't actually in the domain. Then, I'll emphasize that, whenever you solve an equation that involves logs, you should check that the solutions you find are actually in the domain.}

    \begin{problem}[\stretch{0.35}]
      \check Check your solutions by plugging them back into the original equation. 
    \end{problem}

  \end{enumerate}

  

  \pspace{\newpage}

  \begin{enumerate}[resume]
  \item

    \note{If we're short on time, I will jump ahead to the next one.}

    \begin{problem}[\stretch{1.25}]
      $\log_3 (x + 1) - \log_3 x = 2$
    \end{problem}

    \begin{solution}
      We could immediately raise $3$ to the power of both sides,
      but let's use log rules to rewrite the equation first.
      \begin{align*}
        \log_3 (x + 1) - \log_3 x &= 2 \\
        \log_3(\tfrac{x+1}{x})&= 2 \\
        3^{\log_3(\tfrac{x+1}{x})} &= 3^2 \\
        \frac{x+1}{x} &= 9 \\
        x+1 &= 9x \\
        1 &= 8x \\
        x &= \boxed{\frac{1}{8}}
      \end{align*}
    \end{solution}

  \item \label{quadratic-in-log10-x}

    \note{This is an example of an equation where substitution is helpful, so I do want to do this one.}

    \begin{problem}[\vfill]
      $(\log_{10} x)^2 + \log_{10} x = 6$
    \end{problem}

    \begin{solution}
      We can see in this equation that $x$ only appears in the form
      $\log_{10}(x)$, so let's make a substitution, $u=\log_{10}{x}$
      and solve for $u$ first.
      \begin{align*}
        (\log_{10} x)^2 + \log_{10} x = 6 \\
        u^2 + u = 6 \\
        u^2 + u - 6 &= 0 \\
        (u+3)(u-2) &= 0 \\
      \end{align*}
      This means $u=-3$ or $u=2$. Since $u=\log_{10}x$,
      this means that $\log_{10}x =-3$ or $\log_{10}x = 2$.
      This gives us the two solutions $x=\boxed{10^{-3}}$ and 
      $x=\boxed{10^2}$.
    \end{solution}

  \item
    \begin{problem}[\vfill\newpage]
      $\ln \left( 5x^2 \right) = 3 + \ln x$
    \end{problem}

    \begin{solution}
      Let's raise $e$ to both sides.
      \begin{align*}
        e^{\ln(5x^2)} &= e^{3+\ln x} \\
        e^{\ln(5x^2)} &= e^3 \cdot e^{\ln x} \\
        5x^2 &= (e^3)x \\
        5x^2 - (e^3)x &= 0 \\
        x(5x-e^3) &= 0
      \end{align*}
      This gives us the solutions $x=0$ and $x=\tfrac{1}{5}e^3$. However
      $x=0$ is an extraneous solution. So the only solution is
      $x=\boxed{e^3/5}$.
    \end{solution}

  \end{enumerate}

\item 

  \begin{enumerate}

  \item
    \begin{problem}[\vfill]
      $(\log_3 x)^3 = \log_3 x$
    \end{problem}

    \begin{solution}
      Let's make the substitution $u=\log_3 x$.
      \begin{align*}
        (\log_3 x)^3 &= \log_3 x \\
        u^3 &= u \\
        u^3 - u &= 0 \\
        u(u^2 -1) &= 0
      \end{align*}
      This means $u=\log_3 x$ is $-1$, $0$, or $1$. So the solutions
      are $x=\boxed{1/3}$, $x=\boxed{1}$, and $x=\boxed{3}$.
    \end{solution}

  \item
    \begin{problem}[\vfill]
      $2^{2x + 3} = 8^{x - 7}$ (\textit{Hint:} There's exactly one
      solution, and it's an integer.)
    \end{problem}

    \begin{solution}
      First, let's rewrite the equation so that both sides have
      a base of $2$. Then we can take the log base 2 of both sides
      or simply set the exponents equal.
      \begin{align*}
          2^{2x + 3} &= 8^{x - 7} \\
          2^{2x + 3} &= (2^3)^{x - 7} \\
          2^{2x + 3} &= 2^{3(x - 7)} \\
          \log_2(2^{2x + 3}) &= \log_2(2^{3(x - 7)}) \\
          2x+3 &= 3(x-7) \\
          2x + 3 &= 3x - 21 \\
          x &= \boxed{24}.
      \end{align*}
    \end{solution}

  \item
    \begin{problem}[\nstretch{0.75}]
      $\displaystyle 3^x \cdot \frac{5}{3^{x+1}} = 0$
    \end{problem}

    \begin{solution}
      We can start by simplifying the left side to be $\frac{3}{5} = 0$.
      This equation is false, which indicates that there are 
      \fbox{no solutions}.
    \end{solution}

  \item
    \begin{problem}[\vfill]
      $\displaystyle 49^{x} - \frac{\sqrt{7}}{7^{x^2}} = 0$
    \end{problem}

    \begin{solution}
      Let's rewrite the equation so that each side consists of
      one term and then take the log base 7 of both sides.
      \begin{align*}
          49^{x} - \frac{\sqrt{7}}{7^{x^2}} &= 0 \\
          49^x &= \frac{\sqrt{7}}{7^{x^2}} \\
          \log_7(49^x) &= \log_7\!\left(\frac{\sqrt{7}}{7^{x^2}}\right)\\
          \log_7 (49^x) &= \log_7 \sqrt{7} - 
            \log_7 \left(7^{x^2}\right)\\
          x \log_7 49 &= \log_7 \sqrt{7} - x^2 \log_7 7 \\
          2x &= \tfrac{1}{2} - x^2
      \end{align*}
      We can solve this quadratic equation for $x$ using
      the quadratic formula.
      \begin{align*}
          x^2 + 2x - \tfrac{1}{2} &= 0 \\
          x &= \frac{-2 \pm \sqrt{(2)^2 - 4(1)(-\tfrac{1}{2})}}{2(1)} \\
          x &= \boxed{\frac{-2 \pm \sqrt{6}}{2}}
      \end{align*}

      
    \end{solution}

  \item

    \begin{problem}[\vfill]
      $\displaystyle 7 \cdot 4^{x + 2} - \frac{5}{3^{2x + 1}} = 0$
    \end{problem}

    \begin{solution}
      Let's take the natural log of both sides:
      \begin{align*}
        \ln \left( 7 \cdot 4^{x + 2} \right) &= \ln \left(
          \frac{5}{3^{2x + 1}} \right) \\
        \ln 7 + \ln \left( 4^{x + 2} \right) &= \ln 5 - \ln \left( 3^{2x
            + 1} \right) \\
        \ln 7 + (x + 2) \ln 4 &= \ln 5 - (2x + 1) \ln 3 \\
        \intertext{This equation is linear in $x$, so we can solve it by multiplying everything out and then putting all the terms involving $x$ on one side and all the terms not involving $x$ on the other side:}
        \ln 7 + (\ln 4) x + 2 \ln 4 &= \ln 5 - (2 \ln 3) x - \ln 3 \\
        (\ln 4) x + (2 \ln 3) x &= \ln 5 - \ln 3 - \ln 7 - 2 \ln 4 \\
        (\ln 4 + 2 \ln 3) x &= \ln 5 - \ln 3 - \ln 7 - 2 \ln 4 \\
        x &= \boxed{\frac{\ln 5 - \ln 3 - \ln 7 - 2 \ln 4}{\ln 4 + 2 \ln 3}}
      \end{align*}
    \end{solution}

  \end{enumerate}    

\end{enumerate}
\end{document}
